% GENERATED FILE. Edit canonical lessons, then run npm run generate:books.
% canonical-source-hash: fnv1a64:c58516ac6e283f55
% canonical-lessons: IT-C432-alloggiare, IT-C432-pernottare, IT-C432-pernottamento, IT-C432-soggiornare, IT-C432-portiere

\chapter{Pernottare}
\label{ch:it-pernottare}
\hlchaptermodality{432}

\begin{chapteropening}
\textbf{By the end of this chapter:} \emph{I can talk about where I stay: lodging, spending the night, an overnight stay, staying a while and the porter.}

Complete the last lesson of chapter 432: I can talk about where I stay: lodging, spending the night, an overnight stay, staying a while and the porter.
\end{chapteropening}

\section[alloggiare]{alloggiare — to lodge, to stay}

\label{lesson:IT-C432-alloggiare}



\begin{quote}
Before the new one: say the Italian for the cancellation fee, then the Italian for to bring forward.
\end{quote}

\subsection*{You'll want to know: alloggiare}
\textbf{alloggiare} — \textquotedblleft{}to lodge, to stay\textquotedblright{}.

\begin{quote}
Alloggiamo in un piccolo albergo in centro.
\end{quote}

\begin{quote}
\emph{We are staying in a small hotel in the centre.}
\end{quote}

\emph{L'alloggio} — lodging, accommodation.

\begin{cousinweb}
\emph{alloggiare} is built on \emph{loggia} (EtymDB), which was borrowed from Old French \emph{loge} 'shelter, arbour', from Frankish \emph{laubja} 'arbour'.
\end{cousinweb}

\begin{grammarlens}[title={What you've built}]
One more word for telling, explaining and describing at B1.
\end{grammarlens}

\subsection*{Your turn}
\begin{itemize}
\raggedright
  \item \emph{Say it:} \emph{alloggiare}
  \item \emph{Say it:} \emph{alloggiare}, once more
  \item \emph{Say it:} a sentence of your own with \emph{alloggiare}
  \item \emph{Recall:} say the Italian for the cancellation fee, then the Italian for to bring forward, then say \emph{alloggiare} again
\end{itemize}

\begin{tcolorbox}[breakable,colback=teal!4,colframe=teal!35!black,title={Before you move on}]
What does alloggiare mean? (\textquotedblleft{}To lodge, to stay\textquotedblright{}.) Say it once more.
\end{tcolorbox}
\section[pernottare]{pernottare — to stay the night, to spend the night}

\label{lesson:IT-C432-pernottare}



\begin{quote}
Before the new one: say the Italian for to bring forward, then the Italian for to lodge.
\end{quote}

\subsection*{You'll want to know: pernottare}
\textbf{pernottare} — \textquotedblleft{}to stay the night, to spend the night\textquotedblright{}.

\begin{quote}
Pernottiamo a metà strada e ripartiamo all'alba.
\end{quote}

\begin{quote}
\emph{We spend the night halfway and set off again at dawn.}
\end{quote}

Formal; in speech, \emph{passare la notte}.

\begin{cousinweb}
\emph{pernottare} is a learned borrowing from Latin \emph{pernoctāre} 'to stay all night', built from \emph{per} and \emph{nox, noctis} 'night' (Lewis \& Short).
\end{cousinweb}

\begin{grammarlens}[title={What you've built}]
One more word for telling, explaining and describing at B1.
\end{grammarlens}

\subsection*{Your turn}
\begin{itemize}
\raggedright
  \item \emph{Say it:} \emph{pernottare}
  \item \emph{Say it:} \emph{pernottare}, once more
  \item \emph{Say it:} a sentence of your own with \emph{pernottare}
  \item \emph{Recall:} say the Italian for to bring forward, then the Italian for to lodge, then say \emph{pernottare} again
\end{itemize}

\begin{tcolorbox}[breakable,colback=teal!4,colframe=teal!35!black,title={Before you move on}]
What does pernottare mean? (\textquotedblleft{}To stay the night, to spend the night\textquotedblright{}.) Say it once more.
\end{tcolorbox}
\section[il pernottamento]{il pernottamento — the overnight stay}

\label{lesson:IT-C432-pernottamento}



\begin{quote}
Before the new one: say the Italian for to lodge, then the Italian for to stay the night.
\end{quote}

\subsection*{You'll want to know: il pernottamento}
\textbf{il pernottamento} — \textquotedblleft{}the overnight stay\textquotedblright{}.

\begin{quote}
Il prezzo comprende pernottamento e colazione.
\end{quote}

\begin{quote}
\emph{The price includes the overnight stay and breakfast.}
\end{quote}

On a hotel bill, one line per night.

\begin{cousinweb}
\emph{pernottamento} is built on \emph{pernottare}, a learned borrowing from Latin \emph{pernoctāre} 'to stay all night', from \emph{nox} 'night' (Lewis \& Short).
\end{cousinweb}

\begin{grammarlens}[title={What you've built}]
One more word for telling, explaining and describing at B1.
\end{grammarlens}

\subsection*{Your turn}
\begin{itemize}
\raggedright
  \item \emph{Say it:} \emph{il pernottamento}
  \item \emph{Say it:} \emph{il pernottamento}, once more
  \item \emph{Say it:} a sentence of your own with \emph{il pernottamento}
  \item \emph{Recall:} say the Italian for to lodge, then the Italian for to stay the night, then say \emph{il pernottamento} again
\end{itemize}

\begin{tcolorbox}[breakable,colback=teal!4,colframe=teal!35!black,title={Before you move on}]
What does il pernottamento mean? (\textquotedblleft{}The overnight stay\textquotedblright{}.) Say it once more.
\end{tcolorbox}
\section[soggiornare]{soggiornare — to stay (for some time)}

\label{lesson:IT-C432-soggiornare}



\begin{quote}
Before the new one: say the Italian for to stay the night, then the Italian for the overnight stay.
\end{quote}

\subsection*{You'll want to know: soggiornare}
\textbf{soggiornare} — \textquotedblleft{}to stay (for some time)\textquotedblright{}.

\begin{quote}
Abbiamo soggiornato una settimana al lago.
\end{quote}

\begin{quote}
\emph{We stayed a week at the lake.}
\end{quote}

\emph{Il soggiorno} is a stay, and also the living room.

\begin{cousinweb}
\emph{soggiornare} is usually traced through Old French \emph{sojorner} to Vulgar Latin \emph{subdiurnāre} 'to spend the day', built on \emph{diurnus} 'of the day', from \emph{diēs} 'day' (Lewis \& Short).
\end{cousinweb}

\begin{grammarlens}[title={What you've built}]
One more word for telling, explaining and describing at B1.
\end{grammarlens}

\subsection*{Your turn}
\begin{itemize}
\raggedright
  \item \emph{Say it:} \emph{soggiornare}
  \item \emph{Say it:} \emph{soggiornare}, once more
  \item \emph{Say it:} a sentence of your own with \emph{soggiornare}
  \item \emph{Recall:} say the Italian for to stay the night, then the Italian for the overnight stay, then say \emph{soggiornare} again
\end{itemize}

\begin{tcolorbox}[breakable,colback=teal!4,colframe=teal!35!black,title={Before you move on}]
What does soggiornare mean? (\textquotedblleft{}To stay (for some time)\textquotedblright{}.) Say it once more.
\end{tcolorbox}
\section[il portiere]{il portiere — the porter, the doorman; the goalkeeper}

\label{lesson:IT-C432-portiere}



\begin{quote}
Before the new one: say the Italian for the overnight stay, then the Italian for to stay.
\end{quote}

\subsection*{You'll want to know: il portiere}
\textbf{il portiere} — \textquotedblleft{}the porter, the doorman; the goalkeeper\textquotedblright{}.

\begin{quote}
Il portiere ci dà la chiave della camera.
\end{quote}

\begin{quote}
\emph{The porter gives us the room key.}
\end{quote}

In football, \emph{il portiere} is the goalkeeper.

\begin{cousinweb}
\emph{portiere} is built on \emph{porta} 'door', which is inherited from Latin \emph{porta} 'gate' (EtymDB).
\end{cousinweb}

\begin{grammarlens}[title={What you've built}]
One more word for telling, explaining and describing at B1.
\end{grammarlens}

\subsection*{Your turn}
\begin{itemize}
\raggedright
  \item \emph{Say it:} \emph{il portiere}
  \item \emph{Say it:} \emph{il portiere}, once more
  \item \emph{Say it:} a sentence of your own with \emph{il portiere}
  \item \emph{Recall:} say the Italian for the overnight stay, then the Italian for to stay, then say \emph{il portiere} again
\end{itemize}

\begin{tcolorbox}[breakable,colback=teal!4,colframe=teal!35!black,title={Before you move on}]
What does il portiere mean? (\textquotedblleft{}The porter, the doorman; the goalkeeper\textquotedblright{}.) Say it once more.
\end{tcolorbox}
